\silentchapter{3 Seleksjonstrykka produserer eit landskap}{seleksjonstrykka produserer eit landskap}
\prop{3}{}{Seleksjonstrykka produserer eit landskap over formrommet\footnote{Landskapet er summen av alle ytelseskrav og restriksjonar som verkar på systemet.}.}
\prop{3.1}{d}{\textbf{Tilpassingslandskapet\footnote{Ein lokal optimal løysing i eit tenestedesign samsvarar med ein haug i dette abstrakte landskapet.}} er den aggregerte verknaden av alle samtidige seleksjonstrykk over formrommet.\footnote{Wright (1932)}}
\prop{3.11}{t}{Kvar posisjon har ein verdi: kor godt forma stettar alle verksame trykk samstundes.}
\prop{3.111}{o}{Kvart seleksjonstrykk bidreg med ein vekt som varierer over tid.}
\prop{3.1111}{t}{Vektene er ikkje frie parametrar\footnote{Kompilatoren si optimalisering eller marknadens seleksjon verkar som konkrete uttrykk for desse vektene.}.}
\prop{3.11111}{o}{Dei er empirisk tilgjengelege som samvariasjonen mellom kvart trykk og den observerte fordelinga, betinga på alle andre trykk.}
\prop{3.11112}{t}{Landskapet er retrospektivt rekonstruerbart.}
\prop{3.111121}{o}{For eit gjeve tidspunkt kan vektene estimerast frå dei realiserte formene i ein temporal nærleik\footnote{Retro-analyse av ein kodebase let oss utleie nøyaktig kva vekt utviklarane la på yting versus skalerbarheit i det gjevne tidspunktet.}.}
\prop{3.11113}{o}{Landskapet er lokalt prediktivt: det predikerer kva regionar som har høgast sannsynlegheit for neste okkupasjon\footnote{Systemet konvergerer alltid mot den mest stabile arkitekturen i nærleiken av sin noverande tilstand.}.}
\prop{3.111131}{t}{Det predikerer ikkje kva einskildform som vert realisert.}
\prop{3.111132}{o}{Prediksjonshorisonten er ein funksjon av landskapets endringsrate.}
\prop{3.111133}{o}{I periodar med stase er prediksjonen lang; i bifurkasjonar kort.}
\prop{3.12}{t}{Av 2.12, 2.13 og 3.1: tilpassingsfunksjonen har generisk fleire lokale maksima.}
\prop{3.121}{t}{Landskapet har difor fleire haugar.}
\prop{3.1211}{d}{Ein haug er ein posisjon der kvar lita endring gjev dårlegare tilpassing\footnote{Å flytte systemet ut av ein slik tilstand krev eit bevisst energipådrag mot gradienten, for å overkome den suboptimale dalen til neste haug.}.}
\prop{3.12111}{d}{Dalar er posisjonar ingen selekterer.}
\prop{3.121111}{o}{Dei er ikkje forbodne; dei er berre ulønnsame.}
\prop{3.12112}{t}{Eit landskap med éin global optimal er eit spesialtilfelle, ikkje det generiske tilfellet\footnote{Jakta på den universelle, altløysande applikasjonen bryt mot landskapets fundamentale multimodalitet.}.}
\prop{3.121121}{t}{Funksjonalismens implisitte føresetnad om éin optimal løysing er empirisk feil for dei fleste funksjonelle klasser.}
\prop{3.121122}{o}{Multimodalitet i fordelinga er direkte støtte til proposisjon 3.12.}
\prop{3.121123}{o}{Todelinga mellom ein høgrygga og ein lågare modernistisk klynge er ikkje tilfeldig, men eit empirisk mælbart trekk ved tilpassingslandskapet.}
\prop{3.2}{d}{Stabilitetsgraden til ein haug er brattleiken på dalveggane.\footnote{Waddington (1957)\footnote{Kanalisering gjer systemet robust mot brukarfeil og dataavbrot: dalveggane dirigerer feilsteg tilbake til trygg tilstand.}}}
\prop{3.21}{t}{Bratte veggar tyder kanalisering: forma er robust mot forstyrringar.}
\prop{3.211}{o}{Nokre seleksjonstrykk er sterkare enn andre.}
\prop{3.2111}{t}{Dei sterkaste kanaliserer formrommet i få retningar.}
\prop{3.21111}{t}{Hierarkiet av trykk produserer eit hierarki av kanalar: grove kanalar fyrst, finare innanfor dei grove\footnote{Protokollar med høg feil-intoleranse produserer grove, bratte kanalar i formrommet.}.}
\prop{3.211111}{o}{Kanaliseringa er målbar som brattleiken på dalveggane i tettleikskartets topografi.}
\prop{3.211112}{o}{Dei grove kanalane svarer til dimensjonar med låg variasjonskoeffisient; dei fine til dimensjonar med høg variasjonskoeffisient.}
\prop{3.211113}{o}{At nokre dimensjonar er nær låste og andre nær frie er ein direkte empirisk signatur av kanaliseringshierarkiet. Det er ikkje eit artefakt av projeksjonen.}
\prop{3.3}{d}{Ein \textbf{stil} er ei klynge av realiserte former kring same haug.\footnote{Kubler (1962)}}
\prop{3.31}{t}{At ein stil er ei klynge, ikkje ein kategori, forklarer kvifor stilar har uskarpe grenser\footnote{Grensene for kor ein tenestekategori byrjar og ein annan sluttar, kan aldri definerast binært i logikken; overgangen er ein glidande gradient i avstandsmål.}.}
\prop{3.311}{t}{Grensa mellom to klynger i eit kontinuerleg rom er naudsynleg diffus.}
\prop{3.3111}{o}{Stilperiodar er gradientar, ikkje topologiske klynger.}
\prop{3.31111}{o}{Gjennomsnittsstolen i ein stilperiode ligg nærmare nabostilane enn sine eigne medkandidatar.}
\prop{3.31112}{t}{At ein stil er eit statistisk fenomen og ikkje ein ontologisk kategori, forklarer kvifor kantane mellom periodar alltid er akademisk stridomne.}
\prop{3.311121}{t}{Det finst ingen naturleg grense å finne; diffusheita er strukturelt naudsynt.}
\prop{3.4}{t}{Formgjeving konstruerer ikkje landskapet.}
\prop{3.41}{t}{Landskapet er gjeve av vilkåra.}
\prop{3.411}{t}{Å oppdage ein haug er ikkje å skape han\footnote{Når to isolerte utviklarteam konvergerer mot identiske designmønster (som MVC-arkitekturen), provar det at mønsteret er ein stabil haug i landskapet, ikkje ein konvensjon.}.}
\prop{3.4111}{o}{Ei form som dukkar opp uavhengig i tradisjonar utan kontakt, er sterkt prov for at haugen er ein eigenskap ved landskapet, ikkje ved nokon einskild tradisjon.}
\prop{3.41111}{t}{Konvergensen forklarer seg sjølv: same vilkår, same gradient, same haug.}
\prop{3.411111}{t}{Ho krev ingen kommunikasjon mellom tradisjonane; ho krev berre at dei navigerer det same landskapet.}
\prop{3.411112}{i}{At reduktive former oppstår parallelt i tradisjonar utan direkte kontakt, krev ingen intellektuell påverknad som forklaring.}
\prop{3.411113}{t}{Same avgrensingstrykk produserer same haug, uavhengig av kven som navigerer.}

